% from ega0-20-fr.tex
\oldpage[0\textsubscript{IV-1}]{141}
On a donc, en dehors du module d'imperfection $\Upsilon_{C/B/A}$, les
modules d'imperfection $\Upsilon_{B/A/\Lambda}$,
$\Upsilon_{C/A/\Lambda}$ et $\Upsilon_{C/B/\Lambda}$, et on a déjà défini
des homomorphismes canoniques de $C$-modules (20.6.4.2)
\begin{align}
\tag{20.6.14.1}
\label{0.20.6.14.1-fr}
  u':\Upsilon_{C/A/\Lambda}&\longrightarrow\Upsilon_{C/B/\Lambda},\\
\tag{20.6.14.2}
\label{0.20.6.14.2-fr}
  s':\Upsilon_{C/B/\Lambda}&\longrightarrow\Upsilon_{C/B/A}.
\end{align}

Comme dans le diagramme commutatif (20.5.7.3)
\begin{equation}
\tag{20.6.14.3}
\label{0.20.6.14.3-fr}
\xymatrix{
  \Omega^1_{A/\Lambda}\otimes_A B \ar[r] \ar[d] &
  \Omega^1_{B/\Lambda} \ar[r] \ar[d] &
  \Omega^1_{B/A} \ar[r] \ar[d] & 0 \\
  \Omega^1_{A/\Lambda}\otimes_A C \ar[r] &
  \Omega^1_{C/\Lambda} \ar[r] & \Omega^1_{C/A} \ar[r] & 0
}
\end{equation}
la ligne inférieure est formée de $C$-modules, on en déduit par
tensorisation un diagramme commutatif
\begin{equation}
\tag{20.6.14.4}
\label{0.20.6.14.4-fr}
\xymatrix{
  \Omega^1_{A/\Lambda}\otimes_A C \ar[r] \ar@{=}[d] &
  \Omega^1_{B/\Lambda}\otimes_B C \ar[r] \ar[d] &
  \Omega^1_{B/A}\otimes_B C \ar[r] \ar[d] & 0 \\
  \Omega^1_{A/\Lambda}\otimes_A C \ar[r] &
  \Omega^1_{C/\Lambda} \ar[r] & \Omega^1_{C/A} \ar[r] & 0
}
\end{equation}
où la première ligne est encore exacte et la flèche verticale de gauche
est l'identité; si l'on pose
\begin{equation}
\tag{20.6.14.5}
\label{0.20.6.14.5-fr}
  \Upsilon^C_{B/A/\Lambda}=
  \operatorname{Ker}(u_{B/A/\Lambda}\otimes1_C)=
  \operatorname{Ker}(\Omega^1_{A/\Lambda}\otimes_A C\longrightarrow
       \Omega^1_{B/\Lambda}\otimes_B C)
\end{equation}
on voit, compte tenu de la définition de $\Upsilon_{C/A/\Lambda}$, que
l'on a un unique $C$-homomorphisme
\begin{equation}
\tag{20.6.14.6}
\label{0.20.6.14.6-fr}
  v':\Upsilon^C_{B/A/\Lambda}\longrightarrow\Upsilon_{C/A/\Lambda}
\end{equation}
rendant commutatif le diagramme
\begin{equation}
\tag{20.6.14.7}
\label{0.20.6.14.7-fr}
\xymatrix{
  0 \ar[r] & \Upsilon^C_{B/A/\Lambda} \ar[r] \ar[d]^{v'} &
  \Omega^1_{A/\Lambda}\otimes_A C \ar[r] \ar@{=}[d] &
  \Omega^1_{B/\Lambda}\otimes_B C \ar[r] \ar[d] &
  \Omega^1_{B/A}\otimes_B C \ar[r] \ar[d] & 0 \\
  0 \ar[r] & \Upsilon_{C/A/\Lambda} \ar[r] &
  \Omega^1_{A/\Lambda}\otimes_A C \ar[r] &
  \Omega^1_{C/\Lambda} \ar[r] & \Omega^1_{C/A} \ar[r] & 0
}
\end{equation}
dont les lignes sont exactes.

Lorsque $\Lambda=\mathbb{Z}$, on écrira $\Upsilon^C_{B/A}$ au lieu de
$\Upsilon^C_{B/A/\mathbb{Z}}$. Si $\Lambda$ est un corps premier, on a
$\Upsilon^C_{B/A/\Lambda}=\Upsilon^C_{B/A}$.
\end{env}

\begin{env}[20.6.15]
\label{0.20.6.15-fr}
Pour étudier les relations entre les modules précédents, nous introduirons
d'une part le complexe de $B$-modules
$\mathrm{K}_\bullet(B/A/\Lambda)$, d'autre part, les complexes
